\documentclass[10pt]{article}
\usepackage[utf8]{inputenc}
\usepackage[T1]{fontenc}
\usepackage{amsmath}
\usepackage{amsfonts}
\usepackage{amssymb}
\usepackage[version=4]{mhchem}
\usepackage{stmaryrd}
\usepackage{graphicx}
\usepackage[export]{adjustbox}
\graphicspath{ {./images/} }
\usepackage{caption}

\title{MULTI-PARTY PROTOCOLS }

\author{Ashok K. Chandra ${ }^{1}$ Merrick L. Furst ${ }^{2}$ Richard J. Lipton ${ }^{3}$}
\date{}


%New command to display footnote whose markers will always be hidden
\let\svthefootnote\thefootnote
\newcommand\blfootnotetext[1]{%
  \let\thefootnote\relax\footnote{#1}%
  \addtocounter{footnote}{-1}%
  \let\thefootnote\svthefootnote%
}

%Overriding the \footnotetext command to hide the marker if its value is `0`
\let\svfootnotetext\footnotetext
\renewcommand\footnotetext[2][?]{%
  \if\relax#1\relax%
    \ifnum\value{footnote}=0\blfootnotetext{#2}\else\svfootnotetext{#2}\fi%
  \else%
    \if?#1\ifnum\value{footnote}=0\blfootnotetext{#2}\else\svfootnotetext{#2}\fi%
    \else\svfootnotetext[#1]{#2}\fi%
  \fi
}

\begin{document}
\maketitle
\captionsetup{singlelinecheck=false}
by



\begin{enumerate}
  \item Introduction. Many different types of interprocess communication have been examined from a complexity point of view [SP, Y]. We study a new model, in which a collection of processes
\end{enumerate}

$$
P_{0}, \ldots, P_{k-1}
$$

that share information about a set of integers

$$
\left\{a_{0}, \ldots, a_{k-1}\right\},
$$

communicate to determine a 0-1 predicate of the numbers.

For example, suppose $k$ poker players are sitting around a table, and each one is holding a number to his forehead for the others to see. This situation fits our model and we can determine how many bits the players need to communicate in order for all of them to know if the sum of the numbers is greater than $n$.

In this new model, tremendous sharing of information is allowed, while no single party is given enough information to determine the predicate on its own. Formally, each $P_{i}$ has access to every $a_{j}$ except for $a_{i}$.

\footnotetext{${ }^{1}$ Mathematical Sciences Department, IBM Watson Research Center, Yorktown Heights, NY.\\
${ }^{2}$ Department of Computer Science, Carnegie Mellon University, Pittsburgh, Pa.\\
${ }^{3}$ Department of EECS, Princeton University, Princeton, NJ.

Permission to copy without fee all or part of this material is granted provided that the copies are not made or distributed for direct commercial advantage, the ACM copyright notice and the title of the publication and its date appear, and notice is given that copying is by permission of the Association for Computing Machinery. To copy otherwise, or to republish, requires a fee and/or specific permission.
}
For simplicity, we only allow the parties to communicate as follows. At time $t=0$, process $P_{0}$, examining the numbers it knows, broadcasts one bit. At time $t=1$, process $P_{1}$, examining the numbers it knows, and the bit sent by $P_{0}$, broadcasts one bit. This continues in a cyclic fashion until, at a prearranged time , the processes are required to halt, with either all accepting or all rejecting.

The novelty of this shared-information model is apparent once three or more processes are involved. With three processes and three integers, each input can be known by two parties, without any one knowing all three. In this paper we prove both upper and lower bounds on the complexity of multi-party, shared-information protocols (MPP's) for determining certain basic predicates of sets of integers. Our proofs appeal to geometry and Ramsey-like counting arguments. The bounds we prove are unusual in the sense that they are stronger than might naively be expected, and tight-even though we don't know exactly what they are.

Multi-party protocols are a basic and important model of communication in any system, practical or theoretical, where there is shared information. We show that MPP's arise in the study of lower boundsfor branching programs. For this model, the standard\\
approach to proving lower bounds-demonstrating that a certain amount of progress must be made, and that no step makes more than $\delta$ progress, for some small $\delta$-doesn't seem to work. The branching program model is too complicated and there doesn't seem to be a way to define progress appropriately. To prove non-trivial lower bounds a new insight is needed. Our definition of MPP's is motivated by the search for such a new technique. In the last section we make the connection between MPP's and branching programs, and show how the MPP lower bounds can be applied to obtain new lower bounds on time-space trade-offs.\\
2. Preliminaries. From now on we assume that there are $k$-processes $P_{0}, \ldots, P_{k-1}$, and $k$ variables $a_{0}, \ldots, a_{k}$ that range over the integers from 1 to $n$. Each process has access to the $a_{j}$, except that $P_{i}$ is denied access to $a_{i}$. We study what happens when $k$ is fixed and $n$ grows.

Definition. A broadcast history $b \in(0+1)^{*}$ is a record of all the bits transmitted by the processes up to a certain point in time.

DEFINITION. A $k$-party protocol, or just a protocol when $k$ is known, is a deterministic algorithm that, for each process $P_{i}$, determines from the numbers $P_{i}$ knows and the broadcast history, what bit $P_{i}$ should transmit at times $i, i+k, \ldots$.

Definition. Let $H_{n}=\{1, \ldots, n\}^{k}$ be a $k$-dimensional hypercube. Each point $\left\langle v_{0}, v_{1}, \ldots, v_{k-1}\right\rangle \in H_{n_{n}}$ describes a situation, or assignment to the variables $a_{0}, \ldots, a_{k-1}$.

DEFINITION. Let $\Phi$ be a protocol for $k$ processes working with integers in the range 1 to $n$. For each $\bar{v} \in H_{n}, \Phi$ determines a string $\Phi(\bar{v}) \in(0+1)^{*}$ that is the complete broadcast history for the processes in situation $\bar{v}$.

The following technical definition of validity for an assignment of broadcast histories is used to define the more natural notion of validity of protocols, and is motivated by its use in Lemma 2.1.

DEFINITION. Let $Q$ be a $0-1$ predicate of $k$ integers in the range 1 to $n$. An assignment $\Phi$ of broadcast histories to $H_{n}$ is valid for $Q$ if, for every $\bar{v}, \bar{w} \in H_{n}$ that differ in exactly one component,

$$
\Phi(\bar{v}) \neq \Phi(\bar{w}) \text { whenever } Q(\bar{v}) \neq Q(\bar{w}) .
$$

Definition. A protocol $\Phi$ is valid for a predicate $Q$ if $\Phi$ determines a valid broadcast history for $Q$. LEMMA 2.1. A $k$-party protocol $\boldsymbol{\Phi}$ can be used by processes $P_{0}, \ldots, P_{k-1}$ to determine a predicate $Q$ if and only if $\Phi$ is valid for $Q$.

Proof. Suppose $\Phi$ is not valid for $Q$. Then there are two points $\bar{w}$ and $\bar{v}$ such that

\begin{itemize}
  \item[(a)] $Q(\bar{w}) \neq Q(\bar{v})$,
  \item[(b)] $\bar{w}$ and $\bar{v}$ differ only in coordinate $i$, and
  \item[(c)] $\boldsymbol{\Phi}(\bar{w})=\boldsymbol{\Phi}(\bar{v})$.
\end{itemize}

Since process $P_{i}$ cannot see coordinate $i$ it must behave in exactly the same way in situations $\bar{w}$ and $\bar{v}$. Thus, $P_{i}$ must accept or reject incorrectly at one of $\bar{w}$ or $\bar{v}$. Therefore, $\Phi$ cannot be used by the parties to correctly determine $Q$.

The other direction is straightforward. \(\square\)

DEFINITION. The cost of a protocol $\Phi$ on $H_{n}$ is the maximum length of a complete broadcast history on situations in $H_{n}$, i.e.,

$$
\operatorname{cost}_{\Phi}(n)=\max _{\bar{v} \in H_{n}}|\Phi(\bar{v})| .
$$

Definition. The complexity of a predicate $Q$ is the minimum, over all valid protocols $\Phi$ for $Q$, of the cost of $\Phi$. In this setting, complexity is a non-uniform function of $n$.\\[0pt]
3. Lower Bound. We now examine the inherent communication complexity of multi-party protocols. There is a relationship between Ramsey coloring numbers and the number of bits that processes must communicate in order to determine basic predicates. Sec [GRS] or [G] for background on Ramsey theory.

Definition. The predicate Exactly-n is true of $k$ integers $a_{0}, \ldots, a_{k-1}$ if and only if $a_{0}+\cdots+a_{k-1}=n$.

We prove a lower bound on the number of bits $k$ processes have to communicate in order to determine Exactly- $n$. This particular predicate is chosen for convenience, and our results hold for all predicates of this type. As has been noted before, lower bounds on exactly-equal-to predicates imply lower bounds on threshold predicates (" $\sum a_{i} \geq n$ ") [FSS].

The structure of the lower bound proof is as follows. Let $H_{n}=\{1, \ldots, n\}^{k}$ be the set of situations. Let $S_{n}=\left\{\bar{v} \in H_{n} \mid \sum \bar{v}=n\right\}$ be the hyperplane on which the processes must "accept", as shown in Figure 1. We prove that if a protocol determines too few\\
distinct broadcast histories on the plane $S_{n}$, then it is not valid for Exactly- $n$. Thus, to be valid, a protocol must force the parties to transmit enough bits to ensure that there are enough distinct broadcast histories.

\begin{figure}[h]
\begin{center}
  \includegraphics[alt={},max width=\textwidth]{76f4067e-ea33-4374-8907-e685a81b67fa-3_397_664_458_248}
\captionsetup{labelformat=empty}
\caption{Figure 1. $H_{n}$ and $S_{n}$.}
\end{center}
\end{figure}

We first show that protocols do not assign broadcast histories to situations in $H_{n}$ in an arbitrary manner.

Definition. A set of $k$ distinct points $\bar{v}_{0}, \ldots, \bar{v}_{k-1} \in$ $H_{n}$ is a forbidden $k$-pattern if there is a point $\bar{w} \in H_{n}$ such that, for each $i, \bar{w}$ differs from $\bar{v}_{i}$ only in coordinate $i$. (For the case $k=3$, forbidden $k$-patterns on the plane $S_{n}$ are just equilateral triangles.)

LEMMA 3.1. Let $\Phi$ be a protocol for processes $P_{0}$, $\ldots, P_{k-1}$. Let $H_{n}=\{1, \ldots, n\}^{k}$ be the hypercube of situations, and let $\bar{w}=\left\langle w_{0}, \ldots, w_{k-1}\right\rangle$ be any point in $H_{n}$. If $\bar{v}_{0}, \ldots, \bar{v}_{k-1}$ is a forbidden $k$-pattern for $\bar{w}$, and the complete broadcast histories at the $\bar{v}_{j}$ are identical, i.e.,

$$
\alpha=\Phi\left(\bar{v}_{0}\right)=\Phi\left(\bar{v}_{1}\right)=\cdots=\Phi\left(\bar{v}_{k-1}\right),
$$

then $\Phi(\bar{w})=\alpha$.\\
Proof. By induction on the length of the broadcast history we prove that, for all $i$, the $i^{\text {th }}$ bit broadcast at $\bar{w}$ is the same as the $i^{\text {th }}$ bit broadcast at $\bar{v}_{i(\bmod k)}$.\\
Base: If $|\alpha|=0$, then every process immediately halts on $\bar{w}$ since it immediately halts on $\bar{v}_{i}$.

Induction: Suppose, for all histories of length $<t$, the inductive hypothesis is true. Consider the $t^{\text {th }}$ bit to be broadcast. By induction, $P_{t(\bmod k)}$ sees the same broadcast history at $\bar{w}$ and at $\bar{v}_{t(\bmod k)}$, and hence, $P_{t(\bmod k)}$ broadcasts the same bit at time $t$ at $\bar{w}$ as it broadcasts at time $t$ at $\bar{v}_{t(\bmod k)}$.\\
Note that this Lemma is true for all protocols, not just valid ones.

LEMMA 3.2. A $k$-party protocol $\Phi$ is not valid for the predicate Exactly- $n$ if it assigns the same broadcast history to all the points of a forbidden $k$-pattern of $S_{n}$.

Proof. Suppose $\bar{v}_{0}, \ldots, \bar{v}_{k-1}$, is a forbidden $k$-pattern on $S_{n}$ for the point $\bar{w}$, such that

$$
\alpha=\Phi\left(\bar{v}_{0}\right)=\Phi\left(\bar{v}_{1}\right)=\cdots=\Phi\left(\bar{v}_{k-1}\right) .
$$

By geometry, $\bar{w}$ is not on $S_{n}$, and by Lemma 3.1, $\Phi(\bar{w})=\alpha$. Therefore, $\Phi$ is not valid for Exactly- $n$. \(\square\)

DEFINITION. The integer $\chi_{k}(n)$ is the smallest number . of colors required to color the points of $S_{n}$ so that no forbidden $k$-pattern on $S_{n}$ is colored monochromatically.

Theorem 3.3. The complexity of any $k$-party protocol for Exactly-n is bounded below by the logarithm of the number of colors required to color $S_{n}$ so that no forbidden $k$-patterns are colored monochromatically.

Proof. Consider a protocol $\Phi$ for Exactly- $n$. Let each distinct broadcast history $\Phi(\bar{v})$, for $\bar{v} \in H_{n}$, define a color. If some forbidden $k$-pattern on $S_{n}$ is colored monochromatically, then, by Lemma 3.2, $\Phi$ is not valid for Exactly- $n$. This is a contradiction. Therefore, there must be more than $\chi_{k}(n)$ distinct broadcast histories, and hence, in some situation, at least $\log \left(\chi_{k}(n)\right)$ bits must be communicated. \(\square\)

Definition. Let

$$
B_{k}(n)=\left\{\left(v_{0}, \ldots, v_{k-1}\right) \mid 1 \leq v_{i} \leq n\right\} .
$$

For $0 \leq i \leq k-1$, let $\bar{e}_{i}$ be the vector whose $j^{\text {th }}$ coordinate is $\delta_{i j}$.

We now state a powerful Theorem from Ramsey theory. It is actually a corollary of a stronger result and can be thought of as a generalized pigeon-hole principle.

Theorem 3.4. (Gallai [G,GRS]) For every c and k, there exists $a \lambda$ and an $n$, such that every $c$-coloring of $B_{k}(n)$ contains $k$ distinct vectors

$$
\bar{v}, \bar{v}+\lambda \bar{e}_{0}, \ldots, \bar{v}+\lambda \bar{e}_{k-1}
$$

that are identically colored.\\
We are now in a position to prove the main lower bound.

Theorem 3.5. Threshold has multi-party protocol complexity greater than any constant.

Proof. By Theorem 3.3 it suffices to prove that for a fixed $k, \chi_{k}(n)$ is unbounded. To get a contradiction, assume that $\chi_{k}(n)$ is bounded by some constant $c$.

Define the projection $p$ from $S_{n}$ to $B_{k-1}(n)$ by

$$
p\left(x_{0}, \ldots, x_{k-1}\right)=\left(x_{0}, \ldots, x_{k-2}\right) .
$$

Since $\chi_{k}(n) \leq c$, there is a $c$-coloring $C$ of $S_{n}$ in which no forbidden $k$-pattern is colored monochromatically. The mapping $p$ is one-to-one, so $C$ and $p$ induce, in a natural way, a $c$-coloring of the points in $p\left(S_{n}\right)$. Consider a smaller set, $B_{k-1}(n / 4 k)$. It is a subset of $p\left(S_{n}\right)$ and is thus also $c$-colored via $p$. For large enough $n$, Gallai's Theorem implies the existence of $k-1$ distinct vectors

$$
\bar{v}, \bar{v}+\lambda \bar{e}_{0}, \ldots, \bar{v}+\lambda \bar{e}_{k-2},
$$

in $B_{k-1}(n / 4 k)$ that are colored identically.\\
Let $\bar{v}=\left(v_{0}, \ldots, v_{k-2}\right)$, and let

$$
\bar{w}=\left(v_{0}, \ldots, v_{k-2}, n-s-\lambda\right),
$$

where $s=v_{0}+v_{1}+\cdots+v_{k-2}$. Since $s>0$ and $\lambda>0$, the last coordinate of $\bar{w}$ is $\leq n$. Furthermore, $1 \leq n-s-\lambda$ since $s$ is at most $(k-1) n / 4 k$ and $\lambda$ is at most $n / 4 k$. Thus, $\bar{w}$ is a situation in $H_{n}$.

We now have a contradiction since the points

$$
p^{-1}(\bar{v}), \quad p^{-1}\left(\bar{v}+\lambda \bar{e}_{0}\right), \quad \ldots, \quad p^{-1}\left(\bar{v}+\lambda \bar{e}_{k-2}\right)
$$

form a monochromatic forbidden $k$-pattern for $\bar{w}$. \(\square\)\\
4. Upper bound. In this section we prove that, within an additive constant, the lower bound of $\log \left(\chi_{k}(n)\right)$ bits proved in the previous section is also an upper bound. That is, we exhibit a cost $\left[\log \left(\chi_{k}(n)\right)+k\right]$-bit $k$-party protocol that is valid for the predicate Exactly- $n$.

Color the hyperplane $S_{n}=\left\{\bar{v} \in H_{n} \mid \sum \bar{v}=n\right\}$ of $H_{n}$ with $\chi_{k}(n)$ colors such that no forbidden $k$-pattern is monochromatic. Let $\left\langle v_{0}, \ldots, v_{k-1}\right\rangle$ be a situation in $H_{n}$. Let process $P_{i}$ know every number but $v_{i}$. Each process $P_{i}$ computes the point

$$
\bar{w}_{i}=\left\langle v_{0}, \ldots, v_{i-1}, n-\left(\sum_{j \neq i} v_{j}\right), v_{i+1}, \ldots, v_{k-1}\right\rangle \in S_{n} .
$$

At time $t=0, P_{0}$ broadcasts bit 0 of the color at $\bar{w}_{0}$. At $t=1, P_{1}$ broadcasts bit 1 of the color at $\bar{w}_{1}$. This continues until $\log \left(\chi_{k}(n)\right)$ bits are transmitted. Then, each $P_{i}$ in turn transmits a 1 if and only if the color it sees matches the color transmitted in the first $\log \left(\chi_{k}(n)\right)$ bits. Finally, the processes halt and accept if and only if every process braodcast a 1 in the last phase.

THEOREM 4.1. The above protocol is valid for Exactly-n.

Proof. The points $\bar{w}_{i}$ are a forbidden $k$-pattern in $S_{n}$. Thus, if the processes agree that all are the same color, the points must coincide. Hence $\left\langle v_{0}, \ldots, v_{k-1}\right\rangle \in S_{n}$ and the processes correctly accept.

If $\left(v_{0}, \ldots, v_{k-1}\right\rangle \notin S_{n}$, then the processes cannot all see the same color since the forbidden points are not monochromatic. Therefore, the processes correctly reject. \(\square\)

Note that even though $\chi_{k}$ is not known, the optimal protocol can be implemented after an exhaustve search. Although precise bounds are not known, we can relate $\chi_{k}$ to classic Ramsey numbers.

Definition. Let $C_{k}(N)$ be the minimum number of colors needed to color $1, \ldots, N$ such that no length- $k$ arithmetic progression is colored monochromatically.

Theorem 4.2. For $N=k n$,

$$
\chi_{k}(n) \leq C_{k}(N) .
$$

Proof. Define the map $q$ from $S_{n}$ to $\{1, \ldots, N\}$ by

$$
q\left(v_{0}, v_{1}, \ldots, v_{k-1}\right)=v_{0}+2 v_{1}+\cdots+k \cdot v_{k-1} .
$$

Color $1, \ldots, N$ with $c=C_{k}(N)$ colors, avoiding monochromatic length- $k$ arithmetic progressions. Color each $\bar{v} \in S_{n}$ with the color of $q(\bar{v})$. We prove by contradiction that this coloring of $S_{n}$ contains no monochromatic forbidden $k$-patterns.

Assume $\bar{v}_{0}, \ldots, \bar{v}_{k-1}$ is a monochromatic forbidden $k$-pattern. By definition there is a $\bar{w} \in H_{n}$ such that

$$
\bar{v}_{i}=\bar{w}+\lambda_{i} \bar{e}_{i},
$$

for some $\lambda_{0}, \ldots, \lambda_{k-1}$. Since each $\bar{v}_{i}$ is on $S_{n}$, it follows that $\lambda_{0}=\lambda_{1}=\cdots=\lambda_{k-1}$. Consider the points

$$
q\left(\bar{v}_{0}\right), q\left(\bar{v}_{1}\right), \ldots, q\left(\bar{v}_{k-1}\right) .
$$

The map $q$ is linear, so $q\left(\bar{v}_{i}\right)=q(\bar{w})+\lambda_{0} q\left(\bar{e}_{i}\right)$. By definition of $q$,

$$
q\left(\bar{e}_{i}\right)=i+1 ;
$$

hence

$$
q\left(\bar{v}_{i}\right)=a+(i+1) b,
$$

where $a=q(\bar{w})$ and $b=\lambda_{0}$. But then there is a monochromatic arithmetic progression of length $k$, which is a contradiction. \(\square\)

Determining the true rate of growth of $C_{k}(N)$ is a difficult open problem. It is known that $C_{k}$ is unbounded for fixed $k$, however, for $k \geq 4$ it can only be shown to grow extremely slowly. In the remainder of this section we look at the specific case of $k=3$ by relating $C_{k}(N)$ to another Ramsey-like function.

Definition. Let $R_{k}(N)$ be the size of a largest subset $A \subset\{1, \ldots, N\}$ that contains no length- $k$ arithmetic progression.

Theorem 4.3. For some constant $a>0$,

$$
\frac{N}{R_{k}(N)} \leq C_{k}(N) \leq \frac{a N \lg N}{R_{k}(N)} \text {. }
$$

Proof. Suppose $\{1, \ldots, N\}$ can be $c$-colored in such a way that there are no length- $k$ arithmetic progressions. Then some color class must be used at least $N / c$ times. Therefore, $R_{k}(N) \geq N / c$. The first inequality holds since $c=C_{k}(N)$.

Now, assume that $A$ is a subset of $\{1, \ldots, N\}$ that contains no arithmetic progressions of length $k$. We demonstrate that the second inequality holds by proving that it is possible to $c$-color $\{1, \ldots, N\}$ so that it has no length- $k$ arithmetic progressions provided $c \leq a N \lg N /|A|$. This follows directly from:

Lemma 4.4. Let $A \subset\{1, \ldots, N\}$. No more than $O(N \lg N /|A|)$ translates of $A$ are needed to completely cover $\{1, \ldots, N\}$.

Proof of Lemma: We use a probabilistic existence argument. Choose random translations $t_{1}, \ldots, t_{\ell}$, with each $t_{i}$ in the interval $-N$ to $N$. The probability that a value $x \in\{1, \ldots, N\}$ is not covered by a particular translation is at most $1-|A| / 2 N$. Therefore, the expected number of missed values is bounded above by

$$
N \cdot\left(1-\frac{|A|}{2 N}\right)^{\ell} .
$$

For large enough $a$, choosing $\ell=a N \lg N /|A|$ makes this less than 1 and the Lemma is proved. \(\square\)

THEOREM 4.5. (Roth $[R]) R_{3}(n)>n e^{-c \sqrt{\log n}}$. Corollary 4.6. Three players can determine Exactly$n$ in $O(\sqrt{\log n})$ bits.

Proof. From the above,

$$
\chi_{3}(n) \leq C_{3}(3 n) \leq a N \lg N / R_{3}(3 n) \leq a \lg n e^{c \sqrt{\log n}} .
$$

Therefore, $\log \chi_{3}(n)$, the number of bits that have to be communicated, is $O(\sqrt{\log n})$. \(\square\)\\[0pt]
5. Application to Branching Programs. The lower bounds for multi-party protocols give lower bounds on the space and time of branching programs. Branching programs play an important role in the study of time and space tradeoffs and were initially described by Tompa [T] and have been studied by Pippenger $[\mathrm{P}]$ and Borodin, Fischer, Kirkpatrick, Lynch, and Tompa [B]. A lower bound for branching programs of width-2 appears elsewhere in these proceedings [BD].

DEFINITION. A branching program on the bits $b_{0}, \ldots, b_{n-1}$ is a rooted, directed, acyclic graph whose internal nodes are labeled with queries of the form "bit $b_{i}$ ?", whose edges are labeled with possible query answers " $b_{i}=0$ ", or " $b_{i}=1$ ", and whose leaves are labeled either "accept", or "reject". For a given assignment of zeroes and ones to the $b_{i}$, a branching program either accepts or rejects as the path defined by the bit values leads to an accept or reject leaf. Time for a branching program is the maximum path length from the root to a leaf. Space for a branching program is the maximum, over all $d$, of the number of vertices at distance $d$ from the root.

The MPP model described in the first few sections allows us to deal with processes that share information. In the most general situation we can think of processes as having access to a proper subset of a set of bits $b_{0}, \ldots$, $b_{n-1}$. The next lemma shows that the integer-sharing model is also appropriate here.

Lemma 5.1. Let $P_{0}, \ldots, P_{k-1}$ be processes. Let $P_{i}$ have access to subset $s_{i}$ of the bits $\left\{b_{0}, \ldots, b_{n-1}\right\}$. If $\left|s_{i}\right| \leq \alpha \cdot n$, for some $\alpha<1$, i.e., if each process is missing at least a constant fraction of the bits, then there exists a constant $\beta$, and disjoint subsets $r_{0}, \ldots$, $r_{k-1}$ of the $b_{i}$ such that,

\begin{itemize}
  \item[(i)] $s_{i} \cap r_{i}=\emptyset$, i.e., process $P_{i}$ doesn't have access to bits in subset $r_{i}$, and
  \item[(ii)] $\left|r_{i}\right| \geq \beta \cdot n$.
\end{itemize}

Proof. An iterated application of Hall's matching Theorem gives this result directly. \(\square\)

This means that if the bits outside the union of the $r_{i}$ are held constant, the $P_{i}$ share integers whose values are the number of bits "on" in the $r_{j}$.

Theorem 5.2. Any constant-space branching program requires super-linear length to compute Exactly-n.

Proof. Suppose not. Then, for some integers $c, d$ there are branching programs of width $d$ and length $c \cdot n$ that accept Exactly- $n$. Think of a branching program of length $c \cdot n$ as being composed of $2 c d$ equally-long segments. Each segment can be thought of as a process, $P_{i}$, that has access to no more than $n / 2$ bits. The whole branching program can thus be viewed as $2 c d$ processes, each missing $1 / 2$ the bits, communicating a constant amount of information $(2 c d \log d$ bits) and determining Exactly- $n$. Using Lemma 5.1 we can show that this contradicts Theorem 3.3.

Corollary 5.3. Any constant-space branching program for Threshold requires greater than linear time.\\
6. Acknowledgments. The authors wish to thank Lyle McGeoch and Larry Stockmeyer for many helpful suggestions. Special thanks to Bill Scherlis for last minute help with the typesetting. The IBM Thomas Watson Research Center at Yorktown heights provided the creative atmosphere in which this work was begun. This work was partially funded by DARPA contracts N00014082-K-0549 and F33615-81-K-1539.

\section*{References}
\begin{itemize}
  \item[] [B] A. Borodin, M. Fischer, D. Kirkpatrick, N. Lynch, M. Tompa, "A time-space tradeoff for sorting and related non-oblivious computations." Toronto, Dept. of Computer Science, Technical Report 79-01-01, 1979.
  \item[] [BD] A. Borodin, D. Dolev, F. Fich, W. Paul, "Bounds for width-2 branching programs." These proceedings.
  \item[] [EG] P. Erdös and R. Graham, Old and New Problems and Results in Combinatorial Number Theory. L'Enseignement Mathématique, Université de Genève, 1980.
  \item[] [FSS] M. Furst, J. Saxe, M. Sipser, "Parity, circuits and the polynomial-time hierarchy." $22^{\text {nd }}$ Symposium on the Foundations of Computer Science, 1981, pp. 260-270; To appear in Mathematical Systems Theory.
  \item[] [G] R. Graham, Rudiments of Ramsey theory. Regional Conference Series in Mathematics, number 45, 1981.
  \item[] [GRS] R. Graham, B. Rothschild, J. Spencer, Ramsey Theory. Wiley-Interscience, 1980, p. 38.
  \item[] [P] N. Pippenger, "A time-space trade-off." Journal of the ACM, 25, 1978, pp. 509-515.
  \item[] [R] K. Roth, "On certain sets of integers." J. London Math. Soc., 29, 1954, pp. 20-26.
  \item[] [SP] M. Sipser, C. Papadimitriou, "Communication complexity." $14^{\text {th }}$ Annual ACM Symposium on Theory of Computing, 1982, pp. 196-200.
  \item[] [T] M. Tompa, "Time-space tradeoffs for computing functions, using connectivity properties of their circuits." $10^{\text {th }}$ Annual ACM Symposium on Theory of Computing, 1978, pp. 196-204.
  \item[] [Y] A. Yao, "Some complexity questions related to distributed computing." $11^{\text {th }}$ Annual $A C M$ Symposium on Theory of Computing, 1979, pp. 209-213.
\end{itemize}


\end{document}